\newpage
\chapter{Orthogonal Linear Maps}

In this section we are going to take a look at some special type of linear maps that are called \emph{orthogonal} 
and at the term \emph{isometry}. It will be useful to remember some concepts about the dot product such as 
\(\inner{x}{x} = x^T \cdot x = \|x\|^2\).

\subsection{Isometry}

Given two vector spaces \((V, \|\centerdot\|_V)\) and \((W, \|\centerdot\|_W)\). Then a map 
\(f: V \to W\) is called \emph{isometric} if and only if \(\|x - y\|_V = \|f(x) - f(y)\|_W \).

It also has the following traits:

\begin{enumerate}
	\item \emph{Length-Preserving}: \(\|x\|_V = \|y\|_W \quad \forall x \in V\).
	\item \emph{Angle-Preserving}: \(\angle (x,y) = \angle (f(x),f(y)) \quad \forall x,y \in V\)
	\item \emph{Linear}: If \(f(x + \lambda y) = f(x) + \lambda f(y)\)
\end{enumerate}

\subsection{Unit-sphere}

For \((V, \R)\) with metric \(\|\cdot\|\) we denote 

\[
    S = \{x \mid x \in V : \|x\| = 1\}
\]

as the \emph{unit-sphere}.

\subsection{The Rotation Matrix}

This matrix is a common example of an isometric linear map, because of the way it transforms the space.

For the \(2D\) we have 

\[
    \begin{bmatrix}
        \cos(\phi) & -\sin(\phi)\\
        \sin(\phi) & \cos(\phi)
    \end{bmatrix}
\]

The derivation comes from the addition theorems for \(\sin(\phi) \) and \(\cos(\phi)\).

Given two vectors \(\vec{v}\) and \(\vec{g}\) whose angles are \(\alpha\) and \(\theta = \alpha + \beta\). 
We can rotate the vector \(\vec{v}\) by some \(\beta\) to transform into \(\vec{g}\).

\begin{align*}
    x &= r \cos(\alpha)\\
    y &= r \sin(\alpha)
\end{align*}

Then 

\begin{align*}
    x' &= r' \cos(\alpha + \beta)\\
    y' &= r' \sin(\alpha + \beta)
\end{align*}

And therefore

\begin{align*}
    x' &= r' \cos(\alpha)\cos(\beta) - r'\sin(\alpha)\sin(\beta)\\
    y' &= r' \cos(\alpha)\sin(\beta) + r'\cos(\beta)\sin(\alpha)
\end{align*}

Because, a rotation does not change the magnitude we can say that \(r = r'\). Thus 

\begin{align*}
    x' &= r \cos(\alpha)\cos(\beta) - r\sin(\alpha)\sin(\beta)\\
    y' &= r \cos(\alpha)\sin(\beta) + r\cos(\beta)\sin(\alpha)
\end{align*}

Now notice that we have our \(x\) and \(y\) inside the expression

\begin{align*}
    x' &= x\cos(\beta) - y\sin(\beta)\\
    y' &= x\sin(\beta) + y\cos(\beta)
\end{align*}

which gives us 

\[
    \begin{bmatrix}
        \cos(\phi) & -\sin(\phi)\\
        \sin(\phi) & \cos(\phi)
    \end{bmatrix}
    \begin{bmatrix}
        x\\
        y
    \end{bmatrix}
\]

\QED

\subsubsection{Rotation Matrices For \(\R^3\)}

\[
    R_x (\phi) =
    \begin{bmatrix}
        1 & 0 & 0 \\
        0 & \cos(\phi) & -\sin(\phi) \\
        0 & \sin(\phi) & \cos(\phi)
    \end{bmatrix}
\]

\[
    R_y (\phi) =
    \begin{bmatrix}
        1 & 0 & \sin(\phi) \\
        0 & 1 & 0 \\
        -\sin(\phi) & 0 & \cos(\phi)
    \end{bmatrix}
\]

\[
    R_z (\phi) =
    \begin{bmatrix}
        \cos(\phi) & -\sin(\phi) & 0 \\
        \sin(\phi) & \cos(\phi) & 0 \\
        0 & 0 & 1
    \end{bmatrix}
\]

\subsection{Orthogonal Matrices}

A matrix \(A \in \R^{n \times n}\) is \emph{orthogonal} if its columns vectors form an orthonormal 
basis.

The set of all \emph{orthogonal} matrices is \(O(n)\).

\subsubsection{Properties of orthogonal matrices}

\begin{enumerate}
   
    \item \(A \in O(n)\).
   
    \item If \(A\) has an inverse then \(A^{-1} = A^T\).
   
    \item \(A^T \in O(n)\).

    \item \(\|Q\|_2 = \sup_{x \ne 0} \frac{\|Qx\|}{\|x\|_z} = 1\).

\end{enumerate}

\textbf{Proof 1-2:}

\begin{align*}
    A^T = A^{-1} &\iff A^T A = E\\
                 &\iff (A^T A)_{ij} = \inner{a_i}{a_j} = \delta_{ij} 1 \le i,j \le n\\
                 &\iff a_i \bot a_j, i \ne j, \text{ and } \|a_i\| = 1, 1 \le i \le n\\
                 &\iff A \in O(n)
\end{align*}

\textbf{Proof 1-3:}

\begin{align*}
    A \in O(n) &\iff \\
               &\iff A^{-1} = A^T\\
               &\iff (A^T)^{-1} = (A^T)^T \\
               &\iff (A^{1})^{T} = A\\
               &\iff A^T \in O(n)
\end{align*}

\textbf{Example of a proof}

Prove that \(H_n\) is \emph{orthogonal} for all \(\vec{u}\in \R^n \backslash \{0\} \).

\[
    H_n = I_n - 2 \frac{uu^T}{u^Tu}
\]

\(I_n\) is the identity matrix.

\textbf{Proof:}

\begin{align*}
    (H_n)^2 &= I \\
    &= \left(I_n - 2 \frac{uu^T}{u^Tu}\right)^2\\
    &= I_n^2 - 4\frac{uu^T}{u^Tu} + 4\left(\frac{uu^T}{u^Tu}\right)^2\\
    &= I_n -4\frac{uu^T}{u^Tu} +4\frac{uu^T uu^T}{(u^Tu)^2}\\ 
    &= I_n -4\frac{uu^T}{u^Tu} +4\frac{u(u^T u)u^T}{(u^Tu)^2}\\ 
    &= I_n -4\frac{uu^T}{u^Tu} +4\frac{uu^T}{u^Tu}\\ 
    &= I_n
\end{align*}

\QED

\subsection{The Determinants of Orthogonal Matrices}

The \emph{determinant} of this kind of matrices can only be either positioning 1 or negative 1.

\textbf{Proof:}:

We know that \(\det(I) = I \), therefore

\begin{align*}
    \det(I) &= \det(A A^{-1}) \\
            &= \det(A A^T) \\
            &= \det(A) \det(A^T) \\
            &= (\det(A))^2 \\
            &= \|\det (A) \| = 1
\end{align*}

\QED

\subsection{The Orthogonal Group}

\(O(n)\) with the operation matrix multiplication forms the \emph{orthogonal group}.

\subsection{Dot Product}

Given \(f: \R^n \to \R^n\) then this two statements are equivalent.

\[
    \inner{f(x)}{f(y)} = \inner{x}{y} \quad \forall x,y \in \R^n
\]

and 

\(f\) preserves the angles and length.

\(f\) is also called \emph{orthogonal}.

\textbf{Proof \(\Rightarrow\):}
 
\[
    \|x\|^2 = \inner{x}{x} = \inner{f(x)}{f(x)} = \|f(x)\|^2
\]

This implies the length preservation. The isometry gives us:

\begin{align*}
    \|f(x) - f(y)\|^2 
    &= \inner{f(x) - f(y)}{f(x) - f(y)} \\
    &= \inner{f(x)}{f(x)} - 2\inner{f(x)}{f(y)} + \inner{f(y)}{f(y)} \\
    &= \inner{x}{x} - 2\inner{x}{y} + \inner{y}{y} \\
    &= \inner{x - y}{x - y} = \|x - y\|^2
\end{align*}

Then

\begin{align*}
    \angle(f(x), f(y)) 
    &= \arccos\left( \frac{\inner{f(x)}{f(y)}}{\|f(x)\| \cdot \|f(y)\|} \right) \\
    &= \arccos\left( \frac{\inner{x}{y}}{\|x\| \cdot \|y\|} \right) \\
    &= \angle(x, y)
\end{align*}

\QED

\textbf{Proof \(\Leftarrow\):} 

With \(\alpha = \angle(x, y) = \angle(f(x), f(y))\) gilt:

\[
    \inner{f(x)}{f(y)} = \cos(\alpha) \cdot \|f(x)\| \cdot \|f(y)\| = \cos(\alpha) 
    \cdot \|x\| \cdot \|y\| = \inner{x}{y}
\]

\QED

\subsection{Orthonormal Matrices}

A matrix \(Q \in \R^{n \times n}\) is \emph{orthonormal} if for each of its column vectors \(q_j\), they are pairwise orthonormal.

\[
    (q_i, q_j) = \delta_{i,j} = 
    \begin{cases}
        1 & i = j \\ 
        0 & \text{ else }
    \end{cases}
\]

\subsubsection*{Properties of Orthonormal Matrices}

If \(Q, \tilde{Q} \in \R{n \times n}\) are orthonormal matrices, the following applies: 

\begin{enumerate}

    \item \(QQ^T = Q^TQ = I\), i.e. \(Q^{-1} = Q^T\).

    \item \(\kappa_2(Q) = 1\). 

    \item \(Q\tilde{Q}\) is an orthonormal matrix.

\end{enumerate}

\textbf{Proof 1:}

When applying the matrix product, we note that our repeated inner product is applied to each of the 
flipped columns. Since by our definition 

\[
    (q_i, q_j) = \delta_{i,j} = 
    \begin{cases}
        1 & i = j \\ 
        0 & \text{ else }
    \end{cases}
\]

We get a matrix with 1's at the entries \(QQ^t_{i,j}\) if \(i = j\) and else 0. in other words the identity.

\QED

\textbf{Proof 2:}

For all \(x \in \R^n\) holds that 

\[
    \|Qx\|_2^2 = (Qx)^T(Qx) = x^T Q^T Q x = x^T x = \|x\|_2^2
\]

Since orthonormal matrices do not preserve the length. Therefore \(\|Q\|_2 = 1\) and 
analogously, the same for \(Q^{-1}\) hence 

\[
    \kappa_2(Q) = \|Q\|_2 \|Q^{-1}\| = 1
\]

\QED

\textbf{Proof 3:}

Applying our first property, we need to show that \((Q\tilde{Q})^T = (Q\tilde{Q})^{-1}\).

\[
    (Q\tilde{Q})^T(\tilde{Q}Q) = Q^T \tilde{Q}^T (\tilde{Q}Q) = Q^T I Q = Q^T Q = I
\]

Therefore the columns are pairwise orthonormal.

\QED

\subsection{Theorem of the Matrix Representation}

The mapping \( f : \R^n \to \R^n \) is orthogonal if and only if  
the associated matrix representation \( Q \) (with respect to the standard basis) is an orthogonal matrix, i.e., \( Q \in O(n) \).

\textbf{Proof \(\Rightarrow\):} 

\(f\) is linear, and thus has a matrix representation \( Q \). It follows that

\[
    \inner{Qx}{Qy} = \inner{Q^\top Qx}{y} = \inner{x}{y} \quad \forall x, y 
    \in \R^n,
\]

thus,

\[
    \inner{Q^\top Qx - x}{y} = 0 \quad \forall x, y.
\]

Choosing \( y = Q^\top Qx - x \), we get

\[
    \langle Q^\top Qx - x, Q^\top Qx - x \rangle = 0 \quad \forall x,
\]

and hence \( (Q^\top Q - I)x = 0 \quad \forall x \).  
It follows that \( Q^\top Q = I \), and by Theorem 7.3, \( Q \in O(n) \).

\QED

\textbf{Proof \(\Leftarrow\):} 

Follows by the same calculation as above.

\QED

\subsection{Change of Basis}

Let \(f\) be orthogonal. Then the matrix representation of \(f\)  
with respect to any orthonormal basis is orthogonal.

\textbf{Proof:} 

Let \( E \) be the standard basis and \(B\) an orthonormal basis. Let  
\( Q = M_E^E(f) \) and \( S = T_E^B \) be the change-of-basis matrix from \(B\) to \( E \).  
Since the columns of \emph{S} are the basis vectors of \(B\), we have \( S \in O(n) \).  

\[
    M_B^B(f) = S^{-1} Q S \in O(n),
\]

due to the group property of \( O(n) \).

\QED

\subsection{QR-Decomposition}

A use case of orthogonal matrices is the \emph{QR-Decomposition}. The goal is that given a system 
\( Ax = b\) with a matrix \(A\in \R^{m \times n}\), \(m \ge n\) and \(rg(A) = n\) we can find 
two matrices \(Q\) and \(R\) such that \(A = QR \text{ with } R\in \R^{n \times n}\)  
and \(Q \in \R^{m \times n}\). 

\(R\) is an upper triangular matrix and \(Q\) is orthogonal for the columns.

We know that the best approximation by the method of the least squares is given by

\[
    x_s = (A^T A)^{-1} A^T b
\]

Because of the orthogonality of the columns \(rg(Q) = rg(A)\), therefore, it can be inverted. Hence,
we get:

\begin{align*}
    x_s &= (A^T A)^{-1} A^T b \\
        &= ((QR)^T (QR))^{-1} (QR)^t b \\
        &= (R^T Q^T QR)^{-1} (QR)^t b \\
        &= (R^T R)^{-1} R^T Q^T b \\
        &=  R^{-1} (R^T)^{-1} R^T Q^T b \\
        &=  R^{-1} Q^T b \\
    Rx_s&= Q^T b
\end{align*}

Thus, we can solve the linear system by finding \emph{Q} which is matrix with an orthonormal basis of the 
vector space which can be found via the Gram-Schmidt algorithm. The matrix \emph{R} consists of the matrix of the 
inner products our \(A\) matrix and \emph{Q}

\[
    R =
    \begin{bmatrix}
        \inner{q_1}{a_1} & \inner{q_1}{a_2} & \cdots &\inner{q_1}{a_n} \\ 
        0 & \inner{q_2}{a_2} & \cdots & \inner{q_2}{a_n} \\ 
        \vdots & \vdots  & \ddots & \vdots \\
        0 & 0  & \cdots & \inner{q_n}{a_n} \\ 
    \end{bmatrix}
\]

Or even simpler \(R = Q^T A\).

\textbf{Example:}

Find the \emph{QR-Decomposition} of the matrix \(A\) given by:

\[
    \begin{bmatrix}
        1 & 3 \\
        1 & -1 
    \end{bmatrix}
\]

First, we find the orthonormal basis of the columns of our matrix \(A\) 

\[ 
    q_1 = \frac{a_1}{\|a_1\|} = \frac{1}{\sqrt{2}} \begin{bmatrix} 1 \\ 1 \end{bmatrix}
\]

\[ 
    r_2 = a_2 - \inner{a_2}{q_1} q_1 
    = 
    \begin{bmatrix} 1 \\ 3 \end{bmatrix}
    - 
    \langle  
    \begin{bmatrix} 1 \\ 3 \end{bmatrix}
    ,
    \frac{1}{\sqrt{2}} \begin{bmatrix} 1 \\ 1 \end{bmatrix}
    \rangle    
    \frac{1}{\sqrt{2}} 
    \begin{bmatrix} 1 \\ 1 \end{bmatrix}
    = 
    \begin{bmatrix}
        2 \\ -2 
    \end{bmatrix}
\]

\[
    q_2 =
    \frac{1}{\sqrt{8}}
    \begin{bmatrix}
        2 \\ -2 
    \end{bmatrix}
\]

\[
    Q = 
    \begin{bmatrix}
        \frac{1}{\sqrt{2}} & \frac{2}{\sqrt{8}} \\
        \frac{1}{\sqrt{2}} & -\frac{2}{\sqrt{8}}
    \end{bmatrix}
\]

\[
    R =
    \begin{bmatrix}
        \sqrt{2} & \sqrt{2} \\
        0    &    \sqrt{8}
    \end{bmatrix}
\]

